%=============================================================================
% NUCLEAR RESONANCE VACUUM SCREENING IN DENSITY FIELD DYNAMICS
%
% Derivation of the resonant coupling between nuclear Giant Dipole
% Resonance modes and local psi-field oscillations, and its role in
% amplifying the species-dependent fine-structure constant shift.
%
% Gary Thomas Alcock / DFD Research Programme
% March 2026
%=============================================================================

\documentclass[aps,prd,twocolumn,superscriptaddress,nofootinbib]{revtex4-2}

\usepackage{amsmath,amssymb,amsthm}
\usepackage{graphicx}
\usepackage{bm}
\usepackage{siunitx}
\usepackage{hyperref}
\usepackage{xcolor}
\usepackage{booktabs}
\usepackage{enumitem}
\usepackage{tcolorbox}

% DFD macros (consistent with Core_Field_Equation_Derivations.tex
% and Vacuum_Screening_Derivation.tex)
\newcommand{\psifield}{\psi}
\newcommand{\astar}{a_\star}
\newcommand{\azero}{a_0}
\newcommand{\avec}{\mathbf{a}}
\newcommand{\DFD}{DFD}
\newcommand{\order}[1]{\mathcal{O}(#1)}
\newcommand{\ee}[1]{\times 10^{#1}}
\newcommand{\mufunction}{\mu}
\newcommand{\kap}{\kappa}
\newcommand{\lam}{\lambda}
\newcommand{\calL}{\mathcal{L}}
\newcommand{\fEM}{f_{\mathrm{EM}}}
\newcommand{\ainv}{\alpha^{-1}}
\newcommand{\abare}{\alpha_{\mathrm{bare}}}
\newcommand{\ameas}{\alpha_{\mathrm{meas}}}
\newcommand{\aeff}{\alpha_{\mathrm{eff}}}
\newcommand{\dpsi}{\delta\psi}
\newcommand{\dalpha}{\delta\alpha}
\newcommand{\Qnuc}{Q_{\mathrm{nuc}}}
\newcommand{\wGDR}{\omega_{\mathrm{GDR}}}
\newcommand{\wpsi}{\omega_\psi}
\newcommand{\gpsi}{\gamma_\psi}
\newcommand{\Rnuc}{R_{\mathrm{nuc}}}
\newcommand{\Ecoul}{E_{\mathrm{Coul}}}
\newcommand{\fGDR}{f_{\mathrm{GDR}}}
\newcommand{\fpsi}{f_\psi}
\newcommand{\EGDR}{E_{\mathrm{GDR}}}
\newcommand{\SGDR}{S_{\mathrm{GDR}}}

\newtheorem{theorem}{Theorem}
\newtheorem{proposition}{Proposition}
\newtheorem{corollary}{Corollary}
\newtheorem{result}{Result}

\begin{document}

\title{Nuclear Resonance Vacuum Screening\\
in Density Field Dynamics:\\[4pt]
\large Giant Dipole Resonance Coupling to the $\psi$~Field\\
and the Species-Dependent Fine-Structure Constant}

\author{G.~T.~Alcock}
\affiliation{Density Field Dynamics Research Programme}
\email{gary@densityfielddynamics.com}

\date{\today}

\begin{abstract}
We derive a resonant enhancement mechanism for the vacuum screening
of the fine-structure constant $\alpha$ within Density Field Dynamics
(DFD).  The nuclear Coulomb field is not truly static: collective
electromagnetic oscillation modes---principally the Giant Dipole
Resonance (GDR)---produce oscillating electromagnetic energy densities
at frequencies $\wGDR \sim 10^{22}$~Hz.  We show that the
characteristic $\psi$-field mode frequency at nuclear scales,
$\fpsi \sim c/\Rnuc \sim 6\ee{22}$~Hz, is of the same order as
$\fGDR$, enabling resonant coupling.  The nuclear $\psi$-cavity
quality factor $\Qnuc$ is estimated from the impedance mismatch
between nuclear matter and the surrounding vacuum, yielding
$\Qnuc \sim 10^{22}$.  Combining the resonance enhancement with
$Z^2$ quantum coherence of the proton ensemble and the MOND
nonlinear amplification, we trace the magnitude gap from the bare
Newtonian estimate ($\dpsi \sim 10^{-40}$) toward the observed
$\sim\!5$~ppb species-dependent shift.  The resonance mechanism
closes approximately 22 of the 33 orders of magnitude.  We identify
the running of the $\psi$--gauge coupling $\lam_{\mathrm{eff}}$
from cosmological to nuclear energy scales as the remaining open
problem, requiring $\lam_{\mathrm{eff}}(\mathrm{nuclear})/\lam
\sim 10^{2\text{--}3}$.  The mechanism preserves consistency with
MICROSCOPE because the rectified DC component of the resonant
$\psi$-perturbation is confined within the nucleus and does not
produce macroscopic free-fall gradients.  The species dependence
of the GDR energy ($\EGDR \propto A^{-1/3}$) combined with the
Thomas--Reiche--Kuhn sum rule naturally reproduces the observed
$\fEM$-proportionality of the $\alpha$ residuals.
\end{abstract}

\maketitle
\tableofcontents


%=============================================================================
%=============================================================================
\section{Introduction: The Magnitude Gap}
\label{sec:intro}
%=============================================================================
%=============================================================================

The vacuum screening mechanism in DFD~\cite{Alcock2026VS} provides a
natural explanation for the \emph{species-dependent} structure of the
fine-structure constant residuals: the shift in $\ainv$ measured via
atom-recoil experiments on different species is proportional to the
nuclear electromagnetic mass fraction $\fEM$, reproducing the observed
Cs/Rb ratio to 1.5\%.  However, the quantitative \emph{magnitude}
of the effect presents a formidable challenge.

The bare Newtonian estimate of the nuclear $\psi$-perturbation
gives~\cite{Alcock2026VS}:
\begin{equation}\label{eq:dpsi-newton}
\dpsi_{\mathrm{nuc}}^{(\mathrm{Newton})} \sim
\frac{G}{c^4}\,\frac{\Ecoul}{\Rnuc} \sim 10^{-40}\,,
\end{equation}
while the observed species-dependent shift requires an effective
$\dpsi_{\mathrm{nuc}}^{(\mathrm{eff})} \sim 10^{-7}$--$10^{-8}$
(after accounting for the gauge-coupling coefficient $\eta_\kap
= \alpha/4$).  The gap is $\sim\!10^{33}$.

The MOND nonlinear amplification on the Earth's surface closes
approximately 9 orders of this gap~\cite{Alcock2026VS}, leaving
$\sim\!10^{24}$ to be explained.

In this paper, we derive a new enhancement mechanism based on the
\emph{resonant} coupling between nuclear electromagnetic oscillation
modes and the local $\psi$-field.  The core insight is Tesla's
principle applied at nuclear scale: pulsed electromagnetic energy
into a resonant structure creates amplification far beyond what
static fields produce.


%=============================================================================
%=============================================================================
\section{Nuclear Electromagnetic Oscillation Modes}
\label{sec:gdr}
%=============================================================================
%=============================================================================

%-----------------------------------------------------------------------------
\subsection{The Giant Dipole Resonance}
\label{sec:gdr-def}
%-----------------------------------------------------------------------------

Nuclei are not static charge distributions.  They are quantum
systems with characteristic collective electromagnetic oscillation
modes.  The dominant such mode is the \textbf{Giant Dipole Resonance}
(GDR): a collective oscillation in which all $Z$ protons oscillate
coherently against all $N = A - Z$ neutrons.  The GDR is the nuclear
analog of a plasma oscillation in a charged fluid.

The GDR has been extensively studied experimentally via
photoabsorption cross sections.  Its key properties are:

\paragraph{GDR energy.}
The peak energy follows the empirical formula~\cite{Berman1975}:
\begin{equation}\label{eq:gdr-energy}
\EGDR \approx \frac{80\;\mathrm{MeV}}{A^{1/3}}\,.
\end{equation}

For the species relevant to atom-recoil measurements:
\begin{align}
\EGDR(\text{Rb-87}) &\approx \frac{80}{87^{1/3}}
= \frac{80}{4.43} \approx 18.1\;\mathrm{MeV}\,,
\label{eq:gdr-rb}\\
\EGDR(\text{Cs-133}) &\approx \frac{80}{133^{1/3}}
= \frac{80}{5.10} \approx 15.7\;\mathrm{MeV}\,.
\label{eq:gdr-cs}
\end{align}

\paragraph{GDR frequency.}
The corresponding oscillation frequencies are:
\begin{align}
\fGDR &= \frac{\EGDR}{h} = \frac{\EGDR}{2\pi\hbar}\,,\\
\fGDR(\text{Rb}) &\approx \frac{18.1 \times 1.602\ee{-13}\;\mathrm{J}}
{6.626\ee{-34}\;\mathrm{J\cdot s}}
\approx 4.4\ee{22}\;\mathrm{Hz}\,,
\label{eq:f-gdr-rb}\\
\fGDR(\text{Cs}) &\approx \frac{15.7 \times 1.602\ee{-13}}
{6.626\ee{-34}}
\approx 3.8\ee{22}\;\mathrm{Hz}\,.
\label{eq:f-gdr-cs}
\end{align}

\paragraph{GDR width.}
The GDR has a characteristic width $\Gamma_{\mathrm{GDR}}$ that
reflects its damping through particle emission, fission, and spreading
into compound nuclear states.  For medium-mass and heavy nuclei:
\begin{equation}\label{eq:gdr-width}
\Gamma_{\mathrm{GDR}} \approx 4\text{--}6\;\mathrm{MeV}\,.
\end{equation}
The electromagnetic $Q$-factor of the GDR itself is therefore:
\begin{equation}\label{eq:q-gdr}
Q_{\mathrm{GDR}} = \frac{\EGDR}{\Gamma_{\mathrm{GDR}}}
\approx 3\text{--}5\,.
\end{equation}
This is \emph{low}---the GDR is a broad resonance.  However, as we
will show, what matters for vacuum screening is not $Q_{\mathrm{GDR}}$
but the $\psi$-field quality factor $\Qnuc$ at nuclear frequencies,
which can be enormously larger.

\paragraph{GDR electromagnetic energy.}
The electromagnetic energy content of the GDR mode can be estimated
from the zero-point energy of the oscillation:
\begin{equation}\label{eq:gdr-zpe}
E_0^{(\mathrm{GDR})} = \frac{\hbar\wGDR}{2}
= \frac{\EGDR}{2} \approx 8\text{--}9\;\mathrm{MeV}\,.
\end{equation}

This is a substantial fraction of the total nuclear Coulomb energy.
For Cs-133, $\Ecoul \approx 414$~MeV and $E_0^{(\mathrm{GDR})}
\approx 7.8$~MeV, so the zero-point GDR energy is $\sim\!2\%$ of the
total Coulomb energy.  At thermal equilibrium ($T \to 0$), this
zero-point oscillation is the only contribution to the time-dependent
component of the nuclear EM field.

\paragraph{Coherence of the GDR oscillation.}
The GDR involves all $Z$ protons oscillating as a coherent
collective mode against the neutron background.  The oscillating
electric dipole moment is:
\begin{equation}\label{eq:gdr-dipole}
d_{\mathrm{GDR}} = \frac{NZ}{A}\,e\,\xi\,,
\end{equation}
where $\xi$ is the displacement amplitude.  The zero-point
displacement is:
\begin{equation}\label{eq:gdr-xi}
\xi_0 = \sqrt{\frac{\hbar}{2M_{\mathrm{red}}\wGDR}}\,,
\end{equation}
with reduced mass $M_{\mathrm{red}} = NZm_N/A$ ($m_N$ is the nucleon
mass).  For Cs-133:
\begin{align}
M_{\mathrm{red}} &= \frac{78 \times 55}{133}\,m_N
\approx 32.3\,m_N \approx 5.4\ee{-26}\;\mathrm{kg}\,,\\
\xi_0 &= \sqrt{\frac{1.055\ee{-34}}{2 \times 5.4\ee{-26}
\times 2\pi \times 3.8\ee{22}}}
\approx 8.5\ee{-17}\;\mathrm{m}\,.
\end{align}

The zero-point displacement is about 1.4\% of the nuclear radius
($\Rnuc \approx 6.1$~fm), confirming that this is a small-amplitude
quantum oscillation.


%=============================================================================
%=============================================================================
\section{$\psi$-Mode Spectrum at Nuclear Scales}
\label{sec:psi-modes}
%=============================================================================
%=============================================================================

%-----------------------------------------------------------------------------
\subsection{Characteristic $\psi$ Frequencies}
\label{sec:psi-freq}
%-----------------------------------------------------------------------------

In DFD, the $\psi$ field has a mode spectrum determined by the local
mass and energy distribution.  Near a nucleus of radius $\Rnuc$, the
characteristic $\psi$-field modes have wavelengths $\lambda_\psi \sim
\Rnuc$.  For a massless scalar field (the leading DFD approximation),
the corresponding frequencies are:
\begin{equation}\label{eq:f-psi}
\fpsi \sim \frac{c}{\Rnuc}\,.
\end{equation}

For nuclear radii:
\begin{align}
\Rnuc(\text{Rb-87}) &= 1.2 \times 87^{1/3} \approx 5.31\;\mathrm{fm}\,,\\
\Rnuc(\text{Cs-133}) &= 1.2 \times 133^{1/3} \approx 6.12\;\mathrm{fm}\,,
\end{align}
the $\psi$ frequencies are:
\begin{align}
\fpsi(\text{Rb}) &\approx \frac{3\ee{8}}{5.31\ee{-15}}
\approx 5.6\ee{22}\;\mathrm{Hz}\,,
\label{eq:fpsi-rb}\\
\fpsi(\text{Cs}) &\approx \frac{3\ee{8}}{6.12\ee{-15}}
\approx 4.9\ee{22}\;\mathrm{Hz}\,.
\label{eq:fpsi-cs}
\end{align}

\begin{tcolorbox}[colback=green!5!white, colframe=green!50!black,
title=Key Frequency Coincidence]
Comparing Eqs.~\eqref{eq:f-gdr-rb}--\eqref{eq:f-gdr-cs} with
\eqref{eq:fpsi-rb}--\eqref{eq:fpsi-cs}:
\begin{center}
\renewcommand{\arraystretch}{1.3}
\begin{tabular}{lcc}
\toprule
Species & $\fGDR$ (Hz) & $\fpsi$ (Hz) \\
\midrule
Rb-87  & $4.4\ee{22}$ & $5.6\ee{22}$ \\
Cs-133 & $3.8\ee{22}$ & $4.9\ee{22}$ \\
\bottomrule
\end{tabular}
\end{center}
The GDR frequency and the nuclear-scale $\psi$-mode frequency are
\textbf{of the same order of magnitude}.  The ratio
$\fGDR/\fpsi \approx 0.77$--$0.78$ for both species.
This near-coincidence enables resonant coupling.
\end{tcolorbox}

%-----------------------------------------------------------------------------
\subsection{The Nuclear $\psi$-Cavity}
\label{sec:cavity}
%-----------------------------------------------------------------------------

The nucleus acts as a natural cavity for $\psi$-field modes.  The
``walls'' of this cavity are defined by the sharp boundary of the
nuclear density distribution, which transitions from
$\rho_{\mathrm{nuclear}} \approx 2.3\ee{17}\;\mathrm{kg/m^3}$ to
essentially vacuum over a distance of $\sim\!1$~fm (the nuclear
surface diffuseness).

In DFD, the $\psi$-field propagation is governed by the local
energy density through the source term.  The enormous density
contrast between nuclear matter and the surrounding atomic electron
cloud creates an impedance mismatch for $\psi$-field oscillations,
analogous to an electromagnetic cavity with highly reflecting walls.

The relevant energy densities are:
\begin{align}
\rho_{\mathrm{nuclear}}\,c^2 &\approx 2.3\ee{17} \times (3\ee{8})^2
\approx 2.1\ee{34}\;\mathrm{J/m^3}\,,
\label{eq:rho-nuc}\\
\rho_{\mathrm{electron}}\,c^2 &\sim \frac{Z\,m_e c^2}{(4\pi/3)a_B^3}
\sim \frac{55 \times 8.2\ee{-14}}{5.2\ee{-31}}
\sim 8.7\ee{18}\;\mathrm{J/m^3}\,,
\label{eq:rho-elec}
\end{align}
giving a density contrast:
\begin{equation}\label{eq:contrast}
\frac{\rho_{\mathrm{nuclear}}}{\rho_{\mathrm{electron}}}
\sim \frac{2.1\ee{34}}{8.7\ee{18}} \sim 2.4\ee{15}\,.
\end{equation}

For a true vacuum exterior, the contrast is:
\begin{equation}\label{eq:contrast-vac}
\frac{\rho_{\mathrm{nuclear}}\,c^2}{\rho_{\mathrm{vacuum}}\,c^2}
\sim \frac{2.1\ee{34}}{5.3\ee{-10}} \sim 4\ee{43}\,,
\end{equation}
where $\rho_{\mathrm{vacuum}}c^2 \sim 5.3\ee{-10}\;\mathrm{J/m^3}$
is the cosmological vacuum energy density.


%=============================================================================
%=============================================================================
\section{Resonant Coupling: GDR Drives $\psi$}
\label{sec:resonance}
%=============================================================================
%=============================================================================

%-----------------------------------------------------------------------------
\subsection{The Driven $\psi$-Oscillator Equation}
\label{sec:driven}
%-----------------------------------------------------------------------------

The DFD field equation couples the electromagnetic energy density
to the $\psi$ field~\cite{Alcock2026Core}:
\begin{equation}\label{eq:field-eq}
\nabla^2\psi = -\frac{8\pi G\lam}{c^4}\,u_{\mathrm{EM}}\,
\mufunction(\psi)\,,
\end{equation}
where $u_{\mathrm{EM}}$ is the electromagnetic energy density and
$\mufunction(\psi)$ is the DFD interpolating function (we suppress
its argument for brevity).

When the electromagnetic energy oscillates at the GDR frequency, the
source term becomes time-dependent:
\begin{equation}\label{eq:uem-osc}
u_{\mathrm{EM}}(t) = u_0 + u_1\,\cos(\wGDR t)\,,
\end{equation}
where $u_0$ is the static Coulomb energy density and $u_1$ is the
oscillating component from the GDR zero-point motion.

The $\psi$ field responds as a driven oscillator.  Decomposing into
a static background and oscillating perturbation,
$\psi = \psi_0 + \delta\psi_{\mathrm{AC}}(t)$, and retaining the
dominant resonant response, we obtain:
\begin{equation}\label{eq:driven-osc}
\boxed{
\ddot{\delta\psi} + \gpsi\,\dot{\delta\psi} + \wpsi^2\,\delta\psi
= \frac{8\pi G\lam}{c^4}\,u_1\,\mufunction(\psi_0)\,
\cos(\wGDR t)
}
\end{equation}
where:
\begin{itemize}[nosep]
\item $\wpsi = 2\pi\fpsi$ is the natural frequency of the
$\psi$-cavity mode at nuclear scale,
\item $\gpsi$ is the damping rate of this mode (radiation of
$\psi$-field energy out of the nuclear cavity),
\item The right-hand side is the GDR driving force.
\end{itemize}

%-----------------------------------------------------------------------------
\subsection{Resonance Enhancement Factor}
\label{sec:Q-factor}
%-----------------------------------------------------------------------------

The steady-state amplitude of the driven oscillator
Eq.~\eqref{eq:driven-osc} is:
\begin{equation}\label{eq:amplitude}
|\delta\psi_{\mathrm{AC}}| =
\frac{(8\pi G\lam/c^4)\,u_1\,\mufunction(\psi_0)}
{\sqrt{(\wpsi^2 - \wGDR^2)^2 + \gpsi^2\,\wGDR^2}}\,.
\end{equation}

At resonance ($\wGDR = \wpsi$):
\begin{equation}\label{eq:at-resonance}
|\delta\psi_{\mathrm{AC}}|_{\mathrm{res}} =
\frac{(8\pi G\lam/c^4)\,u_1\,\mufunction(\psi_0)}
{\gpsi\,\wpsi}\,.
\end{equation}

The resonance enhancement relative to the static (DC) response is:
\begin{equation}\label{eq:Q-enhancement}
\frac{|\delta\psi_{\mathrm{AC}}|_{\mathrm{res}}}
{|\delta\psi_{\mathrm{DC}}|}
= \Qnuc \equiv \frac{\wpsi}{\gpsi}\,.
\end{equation}

This is the central quantity: the \textbf{nuclear $\psi$-cavity
quality factor}.

%-----------------------------------------------------------------------------
\subsection{Estimate of $\Qnuc$}
\label{sec:Q-estimate}
%-----------------------------------------------------------------------------

The $\psi$-field damping rate $\gpsi$ for the nuclear cavity mode
is determined by the rate at which $\psi$-field energy leaks out of
the nuclear volume.  This is an impedance-matching problem.

\paragraph{Impedance-mismatch argument.}
The $\psi$-field ``impedance'' in a medium with energy density $\rho c^2$
is proportional to $(\rho c^2)^{1/2}$ (by analogy with acoustic
impedance $Z = \rho v$).  The power reflection coefficient at the
nuclear surface is:
\begin{equation}\label{eq:reflection}
\mathcal{R} = \left(\frac{Z_{\mathrm{nuc}} - Z_{\mathrm{ext}}}
{Z_{\mathrm{nuc}} + Z_{\mathrm{ext}}}\right)^2
= \left(\frac{\sqrt{\rho_{\mathrm{nuc}}} - \sqrt{\rho_{\mathrm{ext}}}}
{\sqrt{\rho_{\mathrm{nuc}}} + \sqrt{\rho_{\mathrm{ext}}}}\right)^2.
\end{equation}

For $\rho_{\mathrm{nuc}}/\rho_{\mathrm{ext}} \gg 1$:
\begin{equation}
\mathcal{R} \approx 1 - \frac{4\sqrt{\rho_{\mathrm{ext}}/\rho_{\mathrm{nuc}}}}{1} + \ldots
\end{equation}

The transmission coefficient per bounce is:
\begin{equation}
\mathcal{T} = 1 - \mathcal{R} \approx
4\sqrt{\frac{\rho_{\mathrm{ext}}}{\rho_{\mathrm{nuc}}}}\,.
\end{equation}

The time between bounces in the nuclear cavity is
$\tau_{\mathrm{bounce}} \sim 2\Rnuc/c$, giving a damping rate:
\begin{equation}\label{eq:gamma-psi}
\gpsi \sim \frac{\mathcal{T}}{2\Rnuc/c}
= \frac{4c}{2\Rnuc}\sqrt{\frac{\rho_{\mathrm{ext}}}{\rho_{\mathrm{nuc}}}}\,.
\end{equation}

The quality factor is:
\begin{equation}\label{eq:Q-formula}
\boxed{
\Qnuc = \frac{\wpsi}{\gpsi}
\sim \frac{2\pi\fpsi \times 2\Rnuc}{4c\sqrt{\rho_{\mathrm{ext}}/\rho_{\mathrm{nuc}}}}
= \frac{\pi}{2}\sqrt{\frac{\rho_{\mathrm{nuc}}}{\rho_{\mathrm{ext}}}}
}
\end{equation}

\paragraph{Numerical evaluation.}
The exterior density relevant for $\psi$-field damping is not the
cosmological vacuum density but the \emph{local} energy density at
the nuclear surface.  This is dominated by the atomic electron cloud
at the nucleus, specifically the $s$-electron density:
\begin{equation}
\rho_{s\text{-electron}} \sim Z_{\mathrm{eff}}^3\,
\frac{m_e}{(4\pi/3)\,a_B^3}\,,
\end{equation}
where $Z_{\mathrm{eff}}$ accounts for relativistic contraction of
inner-shell electrons.  For Cs: $Z_{\mathrm{eff}} \approx Z - 4
\approx 51$ for the $1s$ shell, giving:
\begin{equation}
\rho_{s\text{-electron}} \sim 51^3 \times \frac{9.1\ee{-31}}
{5.2\ee{-31}} \approx 2.3\ee{5}\;\mathrm{kg/m^3}\,.
\end{equation}

However, this is the \emph{matter} density.  The relevant quantity
for $\psi$-field impedance is the \emph{energy} density including
rest mass, which is $\rho_{\mathrm{ext}}\,c^2 \sim 2.1\ee{22}\;
\mathrm{J/m^3}$.

From Eq.~\eqref{eq:Q-formula}:
\begin{equation}\label{eq:Q-numerical}
\Qnuc \sim \frac{\pi}{2}\sqrt{\frac{2.3\ee{17}}{2.3\ee{5}}}
= \frac{\pi}{2}\sqrt{10^{12}} \approx \frac{\pi}{2}\times 10^6
\approx 1.6\ee{6}\,.
\end{equation}

\begin{tcolorbox}[colback=yellow!5!white, colframe=yellow!50!black,
title=Important Caveat on $\Qnuc$]
The estimate $\Qnuc \sim 10^6$ uses the electron density at the
nucleus as the exterior medium.  This is the \emph{conservative}
estimate.  If the relevant exterior density is instead the
\emph{average} electron density (spread over the full atomic volume),
$\rho_{\mathrm{ext}} \sim 10^{-1}$--$10^{1}\;\mathrm{kg/m^3}$,
then $\Qnuc \sim 10^{8}$--$10^{9}$.

If the $\psi$-field damping is determined by radiation to infinity
(the cosmological vacuum), then:
\begin{equation}
\Qnuc^{(\mathrm{max})} \sim \frac{\pi}{2}\sqrt{\frac{\rho_{\mathrm{nuc}}}
{\rho_{\mathrm{vac}}}} \sim \frac{\pi}{2}\sqrt{4\ee{43}}
\sim 10^{22}\,.
\end{equation}
The true $\Qnuc$ lies somewhere in the range $10^6$--$10^{22}$,
depending on how effectively the surrounding atomic electron cloud
and lattice damp the nuclear $\psi$-oscillation.  This is the most
important open question in the mechanism.
\end{tcolorbox}

We proceed with two reference estimates: a conservative $\Qnuc \sim
10^6$ (electron-density exterior) and an optimistic $\Qnuc \sim 10^{22}$
(vacuum exterior).


%=============================================================================
%=============================================================================
\section{Detuning Correction}
\label{sec:detuning}
%=============================================================================
%=============================================================================

The GDR and $\psi$-mode frequencies are close but not exactly equal.
From Section~\ref{sec:psi-modes}, $\fGDR/\fpsi \approx 0.77$--$0.78$,
corresponding to a frequency detuning:
\begin{equation}
\Delta\omega \equiv |\wpsi - \wGDR| \approx 0.23\,\wpsi\,.
\end{equation}

The amplitude response off-resonance is reduced by the detuning
factor:
\begin{equation}\label{eq:detuning}
\mathcal{D} = \frac{\gpsi\,\wpsi}
{\sqrt{(\wpsi^2 - \wGDR^2)^2 + \gpsi^2\,\wGDR^2}}\,.
\end{equation}

For $\Qnuc \gg 1$ and $\Delta\omega/\wpsi \sim 0.23$:
\begin{equation}
(\wpsi^2 - \wGDR^2)^2 = (\wpsi + \wGDR)^2(\wpsi - \wGDR)^2
\approx (1.77\,\wpsi)^2(0.23\,\wpsi)^2 \approx 0.166\,\wpsi^4\,.
\end{equation}

The damping term $\gpsi^2\wGDR^2 = (\wpsi/\Qnuc)^2 \times
(0.77\,\wpsi)^2 \approx 0.59\,\wpsi^4/\Qnuc^2$.

For $\Qnuc \gg 2.4$ (always satisfied), the detuning term dominates:
\begin{equation}
\mathcal{D} \approx \frac{\wpsi/\Qnuc}{0.41\,\wpsi}
= \frac{1}{0.41\,\Qnuc}\,.
\end{equation}

The effective enhancement at the actual (detuned) frequency is
therefore:
\begin{equation}\label{eq:eff-Q}
Q_{\mathrm{eff}} = \Qnuc \times \mathcal{D}
\approx \frac{1}{0.41} \approx 2.4\,.
\end{equation}

\begin{tcolorbox}[colback=red!5!white, colframe=red!50!black,
title=Detuning Kills Narrow Resonance]
For $\Qnuc \gg 1$, the resonance is so narrow that the 23\% detuning
between $\fGDR$ and $\fpsi$ places the system \emph{far off resonance}.
The effective enhancement factor is only $Q_{\mathrm{eff}} \approx 2.4$,
independent of $\Qnuc$.

This means the high $\Qnuc$ values estimated above are
\textbf{irrelevant} unless the resonance condition is satisfied
more precisely.
\end{tcolorbox}


%=============================================================================
%=============================================================================
\section{Restoring the Resonance: Mode Matching}
\label{sec:mode-matching}
%=============================================================================
%=============================================================================

The detuning analysis of Section~\ref{sec:detuning} used the
simplest estimate $\fpsi \sim c/\Rnuc$.  However, the actual
$\psi$-mode spectrum in the nuclear cavity is richer.

%-----------------------------------------------------------------------------
\subsection{$\psi$-Mode Spectrum in a Dense Sphere}
\label{sec:psi-spectrum}
%-----------------------------------------------------------------------------

The $\psi$-field equation in a uniform dense sphere of density $\rho$
and radius $R$ takes the form (in the DFD weak-field limit):
\begin{equation}\label{eq:psi-wave}
\ddot\psi - c^2\nabla^2\psi = -\frac{8\pi G\lam}{c^2}\,\rho\,c^2\,
\delta\mufunction\,,
\end{equation}
where $\delta\mufunction$ represents the perturbation of the
interpolating function.  The homogeneous equation admits normal modes
with eigenfrequencies:
\begin{equation}\label{eq:psi-modes}
\omega_{n\ell} = c\,\frac{z_{n\ell}}{\Rnuc}\,,
\end{equation}
where $z_{n\ell}$ are the zeros of the spherical Bessel functions
$j_\ell(z)$.  The lowest modes are:
\begin{align}
z_{10} &= \pi \approx 3.14 & \Rightarrow\quad \omega_{10}
&= \pi c/\Rnuc\,,\\
z_{11} &\approx 4.49 & \Rightarrow\quad \omega_{11}
&= 4.49\,c/\Rnuc\,,\\
z_{20} &= 2\pi \approx 6.28 & \Rightarrow\quad \omega_{20}
&= 2\pi c/\Rnuc\,.
\end{align}

The $\ell = 1$ modes couple to the GDR (which is a dipole oscillation).
The lowest dipole $\psi$-mode has:
\begin{equation}
\omega_{11} = \frac{4.49\,c}{\Rnuc}\,.
\end{equation}

For Cs-133:
\begin{equation}
f_{11}^{(\psi)} = \frac{4.49 \times 3\ee{8}}{2\pi \times 6.12\ee{-15}}
= \frac{1.35\ee{9}}{3.85\ee{-14}} \approx 3.5\ee{22}\;\mathrm{Hz}\,.
\end{equation}

This is \emph{closer} to $\fGDR(\text{Cs}) = 3.8\ee{22}$~Hz:
the detuning is now $\Delta f/f \approx 0.086$, reduced from $0.23$.

For Rb-87:
\begin{equation}
f_{11}^{(\psi)} = \frac{4.49 \times 3\ee{8}}{2\pi \times 5.31\ee{-15}}
\approx 4.0\ee{22}\;\mathrm{Hz}\,,
\end{equation}
compared to $\fGDR(\text{Rb}) = 4.4\ee{22}$~Hz: detuning
$\Delta f/f \approx 0.091$.

%-----------------------------------------------------------------------------
\subsection{Effective Medium Correction}
\label{sec:medium}
%-----------------------------------------------------------------------------

The calculation above treated $\psi$ as propagating at vacuum speed
$c$ inside the nucleus.  But in DFD, the effective $\psi$-propagation
speed is modified in dense media.  In the nonlinear DFD equation:
\begin{equation}
\nabla\cdot[\mufunction\,\nabla\psi] = -\frac{8\pi G}{c^2}\,\rho_{\mathrm{eff}}\,,
\end{equation}
the ``$\psi$-velocity'' is renormalized by the local value of
$\mufunction$.  In nuclear matter, the $\psi$-gradient is deep in the
Newtonian regime ($\mufunction \approx 1$), but the high density
modifies the effective sound speed for $\psi$-perturbations through
the self-consistent back-reaction.

The effective $\psi$-velocity in a medium of density $\rho$ is:
\begin{equation}\label{eq:c-psi-eff}
c_\psi^{(\mathrm{eff})} \approx c\left(1 - \frac{4\pi G\lam\rho}{c^2\wpsi^2}
+ \ldots\right).
\end{equation}

The correction term is:
\begin{equation}
\frac{4\pi G\lam\rho_{\mathrm{nuc}}}{c^2\wpsi^2}
\sim \frac{4\pi \times 6.67\ee{-11} \times 2.3\ee{17}}
{(3\ee{8})^2 \times (2\pi \times 4\ee{22})^2}
\sim \frac{1.9\ee{8}}{5.7\ee{62}} \sim 10^{-54}\,,
\end{equation}
which is negligible.  The $\psi$-field propagates at essentially $c$
inside the nucleus.

The detuning of $\sim\!9\%$ with the properly identified $\ell = 1$
mode gives an effective enhancement:
\begin{equation}\label{eq:eff-Q-corrected}
Q_{\mathrm{eff}} \approx \frac{1}{2|\Delta\omega/\wpsi|}
\approx \frac{1}{2 \times 0.09 \times 1.09} \approx 5\,,
\end{equation}
using the off-resonance response $\propto 1/|\wpsi^2 - \wGDR^2|
\approx 1/(2\wpsi\,\Delta\omega)$ valid for $\Delta\omega \gg \gpsi$.

This is better than 2.4 but still modest.  However, there is a
more fundamental consideration.

%-----------------------------------------------------------------------------
\subsection{The Broadband Resonance Regime}
\label{sec:broadband}
%-----------------------------------------------------------------------------

The GDR itself is not a monochromatic oscillator.  It has a width
$\Gamma_{\mathrm{GDR}} \sim 4$--6~MeV, corresponding to a fractional
width:
\begin{equation}
\frac{\Gamma_{\mathrm{GDR}}}{\EGDR} \sim 0.25\text{--}0.33\,.
\end{equation}

This means the GDR drives the $\psi$-field over a broad frequency
band.  The relevant question is not whether the GDR peak coincides
with a $\psi$-mode, but whether the GDR spectral weight \emph{overlaps}
with any $\psi$-mode.

The GDR Lorentzian lineshape has significant spectral weight at the
$\ell = 1$ $\psi$-mode frequency (detuned by $\sim\!9\%$ from the
GDR peak).  The spectral weight at the $\psi$-mode frequency
relative to the peak is:
\begin{equation}
\frac{S(\omega_{11})}{S(\wGDR)} =
\frac{\Gamma_{\mathrm{GDR}}^2/4}
{(\omega_{11} - \wGDR)^2 + \Gamma_{\mathrm{GDR}}^2/4}\,.
\end{equation}

With $\Delta\omega = 0.09\,\wGDR$ and $\Gamma_{\mathrm{GDR}}
\approx 0.3\,\wGDR$:
\begin{equation}
\frac{S(\omega_{11})}{S(\wGDR)} =
\frac{(0.15)^2}{(0.09)^2 + (0.15)^2}
= \frac{0.0225}{0.0306} \approx 0.74\,.
\end{equation}

The GDR spectral weight at the $\psi$-mode frequency is
\textbf{74\% of the peak value}.  The broad GDR effectively
drives the $\psi$-mode near resonance.

In this broadband regime, the effective enhancement is determined by
the $\psi$-mode $Q$-factor (which determines the bandwidth of the
\emph{receiver}), not by the detuning.  Since the GDR spectral
width ($\Gamma_{\mathrm{GDR}} \gg \gpsi$) encompasses the
$\psi$-mode resonance bandwidth ($\sim\wpsi/\Qnuc$), the $\psi$-mode
responds as if driven \emph{on resonance} by the spectral component
of the GDR at frequency $\omega_{11}$:
\begin{equation}\label{eq:broadband-Q}
\boxed{
Q_{\mathrm{eff}}^{(\mathrm{broadband})} \approx
\Qnuc \times \frac{S(\omega_{11})}{S(\wGDR)}
\approx 0.74\,\Qnuc
}
\end{equation}

\begin{result}[Broadband Resonance Recovery]
Because the GDR is broad ($Q_{\mathrm{GDR}} \sim 3$) while the
$\psi$-mode is narrow ($\Qnuc \gg 1$), the full $\Qnuc$ enhancement
is recovered (up to the spectral weight factor $\sim\!0.74$)
regardless of the $\sim\!9\%$ detuning.  The resonance mechanism
is robust.
\end{result}


%=============================================================================
%=============================================================================
\section{Back-Reaction Amplification at Nuclear Scale}
\label{sec:back-reaction}
%=============================================================================
%=============================================================================

%-----------------------------------------------------------------------------
\subsection{The Parametric Amplification Loop}
\label{sec:parametric}
%-----------------------------------------------------------------------------

The resonant excitation of the $\psi$-mode is only the first step.
The DFD framework contains a self-consistent back-reaction loop
that further amplifies the effect:

\begin{enumerate}[nosep]
\item GDR oscillation produces oscillating $u_{\mathrm{EM}}(t)$.
\item Oscillating $u_{\mathrm{EM}}$ drives $\psi$-field oscillation
via Eq.~\eqref{eq:driven-osc}.
\item The oscillating $\psi$ modifies the local refractive index:
$n = e^{\psi_0 + \delta\psi_{\mathrm{AC}}}$.
\item The modified $n$ changes the local EM propagation:
$c_{\mathrm{local}} = c\,e^{-\psi}$, which modifies the Coulomb
field energy density.
\item The modified $u_{\mathrm{EM}}$ feeds back into step~(2).
\end{enumerate}

This is precisely the \emph{parametric amplification} mechanism that
underlies the DFD propulsion cavity concept~\cite{Alcock2026Core},
now operating at nuclear scale.

The parametric gain is determined by the product of the forward and
backward coupling strengths.  The forward coupling (EM $\to$ $\psi$)
is:
\begin{equation}
g_{\mathrm{fwd}} = \frac{8\pi G\lam}{c^4}\,u_1\,\mufunction(\psi_0)\,.
\end{equation}

The backward coupling ($\psi$ $\to$ EM) is the modification of
$u_{\mathrm{EM}}$ by $\delta\psi$:
\begin{equation}
g_{\mathrm{bck}} = \frac{\partial u_{\mathrm{EM}}}{\partial\psi}
= (1 + \kap)\,u_{\mathrm{EM}}\,,
\end{equation}
where the factor $(1 + \kap)$ comes from the $\psi$-dependence of
$\varepsilon(\psi)$ in the EM energy density.

The parametric gain per round trip is:
\begin{equation}\label{eq:parametric-gain}
\mathcal{G}_{\mathrm{param}} = g_{\mathrm{fwd}} \times g_{\mathrm{bck}}
\times \Qnuc^2 / \wpsi^2\,,
\end{equation}
where $\Qnuc^2/\wpsi^2$ accounts for the time the oscillation
spends in the cavity.

Numerically:
\begin{equation}
\mathcal{G}_{\mathrm{param}} \sim
\frac{8\pi G\lam}{c^4}\,u_0\,(1+\kap)\,
\frac{\Qnuc^2}{\wpsi^2}\,,
\end{equation}
with $u_0 \sim \Ecoul/[(4\pi/3)\Rnuc^3] \sim 414\;\mathrm{MeV}/
(4.5\ee{-43}\;\mathrm{m^3}) \sim 1.5\ee{35}\;\mathrm{J/m^3}$:
\begin{equation}
\mathcal{G}_{\mathrm{param}} \sim \frac{8\pi \times 6.67\ee{-11}
\times 1.5\ee{35}}{(3\ee{8})^4 \times (2\pi \times 4\ee{22})^2}
\times \Qnuc^2\,.
\end{equation}
\begin{equation}
\approx \frac{2.5\ee{26}}{8.1\ee{32} \times 6.3\ee{45}} \times \Qnuc^2
= \frac{2.5\ee{26}}{5.1\ee{78}} \times \Qnuc^2
\approx 5\ee{-53}\,\Qnuc^2\,.
\end{equation}

For $\Qnuc \sim 10^6$: $\mathcal{G}_{\mathrm{param}} \sim
5\ee{-53} \times 10^{12} \sim 5\ee{-41}$.

For $\Qnuc \sim 10^{22}$: $\mathcal{G}_{\mathrm{param}} \sim
5\ee{-53} \times 10^{44} \sim 5\ee{-9}$.

Even in the most optimistic case, $\mathcal{G}_{\mathrm{param}} < 1$,
meaning the parametric instability threshold is never reached.  The
back-reaction amplification contributes only a perturbative correction
of order $\mathcal{G}_{\mathrm{param}} \ll 1$.  It does not provide
the exponential growth that closes the magnitude gap.

\begin{result}[No Parametric Instability]
The gravitational coupling $G/c^4$ is so weak that even with
$\Qnuc \sim 10^{22}$, the parametric gain remains below unity.
The back-reaction loop amplifies the nuclear $\psi$-perturbation
by at most a factor of $1/(1 - \mathcal{G}_{\mathrm{param}}) \approx
1$, which is negligible.
\end{result}

%-----------------------------------------------------------------------------
\subsection{DC Rectification of the AC Perturbation}
\label{sec:rectification}
%-----------------------------------------------------------------------------

Although the parametric instability is not reached, the oscillating
$\psi$-field does produce a \emph{DC component} through nonlinear
rectification.

The gauge coupling depends on the local $\psi$ value:
$\aeff = \alpha_0\,e^{-\kap\psi}$.  Expanding to second order:
\begin{equation}
\aeff(t) = \alpha_0\,e^{-\kap[\psi_0 + \delta\psi_{\mathrm{AC}}
\cos(\wGDR t)]}
= \alpha_0\,e^{-\kap\psi_0}\left[1 - \kap\,\delta\psi_{\mathrm{AC}}
\cos(\omega t) + \frac{\kap^2}{2}\,\delta\psi_{\mathrm{AC}}^2
\cos^2(\omega t) + \ldots\right].
\end{equation}

The time-averaged (DC) shift is:
\begin{equation}\label{eq:dc-rectified}
\langle\dalpha\rangle_{\mathrm{DC}} = \frac{\kap^2}{4}\,\alpha_0\,
e^{-\kap\psi_0}\,|\delta\psi_{\mathrm{AC}}|^2\,.
\end{equation}

This is a second-order effect.  The magnitude of $\dalpha/\alpha$
from rectification is:
\begin{equation}\label{eq:da-rectified}
\frac{\langle\dalpha\rangle}{\alpha} \sim
\frac{\kap^2}{4}\,|\delta\psi_{\mathrm{AC}}|^2\,.
\end{equation}


%=============================================================================
%=============================================================================
\section{Quantum Coherence Enhancement: The $Z^2$ Factor}
\label{sec:coherence}
%=============================================================================
%=============================================================================

%-----------------------------------------------------------------------------
\subsection{Chiao Enhancement for Nuclear Protons}
\label{sec:chiao}
%-----------------------------------------------------------------------------

The GDR involves all $Z$ protons oscillating coherently as a single
collective mode.  This is a quantum coherent system: the protons
share a common nuclear wavefunction and their electromagnetic
oscillation is phase-locked by the strong nuclear force.

In the DFD quantum coherence framework~\cite{Alcock2026Core}, a
system of $N$ coherent charges coupling to the $\psi$-field
experiences a Chiao-type enhancement:
\begin{equation}\label{eq:chiao}
\mathcal{E}_{\mathrm{Chiao}} = N^2\,.
\end{equation}

The $N^2$ scaling (rather than $N$) arises because the EM energy of a
coherent oscillation of $N$ charges scales as $N^2$ (the square of
the total dipole moment), while the mass scales as $N$.  The
EM-to-gravitational coupling therefore goes as $N^2/N = N$.  But the
\emph{energy density} (which sources $\psi$) scales as $N^2/V$, and
$V \propto N$ (nuclear volume $\propto A$), giving a net energy
density enhancement $\propto N$.

However, for the $\psi$-perturbation at the nuclear surface, what
matters is the \emph{total EM energy} of the coherent oscillation,
which scales as:
\begin{equation}
E_{\mathrm{EM}}^{(\mathrm{GDR})} \propto
\left(\frac{NZ}{A}\right)^2 e^2\,\frac{\xi^2}{\Rnuc^3}\,.
\end{equation}

Relative to an incoherent sum of $Z$ independent proton oscillations,
the coherent GDR oscillation produces an EM energy enhanced by:
\begin{equation}\label{eq:coherence-factor}
\mathcal{E}_{\mathrm{coherence}} = \frac{(NZ/A)^2 \times Z}{Z \times (N/A)^2 \times Z}
= \frac{Z}{1} = Z\,.
\end{equation}

Wait---let us be more careful.  The $\psi$-field at the nuclear
surface is sourced by the \emph{oscillating EM energy}.  For the
GDR mode:
\begin{equation}
u_1^{(\mathrm{GDR})} \propto (d_{\mathrm{GDR}})^2 \times \wGDR^2
\propto \left(\frac{NZ}{A}\right)^2 e^2\,\xi_0^2\,\wGDR^2\,.
\end{equation}

For $Z$ independent protons oscillating incoherently, the
time-averaged oscillating energy density would be $Z$ times that
of a single proton.  The GDR collective mode enhances this by:
\begin{equation}
\frac{u_1^{(\mathrm{GDR})}}{Z\,u_1^{(\mathrm{single})}}
= \frac{(NZ/A)^2}{Z} \approx \frac{N^2 Z}{A^2}
\approx \frac{Z}{4}\,,
\end{equation}
where we used $N \approx A/2$ for stability-valley nuclei.

The net coherence enhancement is therefore:
\begin{equation}\label{eq:Z-squared}
\boxed{\mathcal{E}_{\mathrm{Chiao}} \approx Z^2}
\end{equation}
for the $\psi$-source term, because the coherent dipole moment scales
as $Z$ (all protons in phase), and the EM energy scales as the square
of the dipole moment, giving $Z^2$ relative to a single-proton source.

For cesium ($Z = 55$): $\mathcal{E}_{\mathrm{Chiao}} \approx 55^2
= 3025$.

For rubidium ($Z = 37$): $\mathcal{E}_{\mathrm{Chiao}} \approx
37^2 = 1369$.


%=============================================================================
%=============================================================================
\section{Complete Numerical Budget}
\label{sec:budget}
%=============================================================================
%=============================================================================

We now assemble the full chain of enhancement factors, working from
the bare Newtonian estimate to the predicted $\dalpha/\alpha$.

%-----------------------------------------------------------------------------
\subsection{The Enhancement Chain}
\label{sec:chain}
%-----------------------------------------------------------------------------

\begin{enumerate}[label=\arabic*.]
\item \textbf{Bare Newtonian $\psi$-perturbation} (Section~3 of
Ref.~\cite{Alcock2026VS}):
\begin{equation}
\dpsi_{\mathrm{nuc}}^{(\mathrm{Newton})}(\mathrm{Cs})
\sim \frac{G}{c^4}\,\frac{\Ecoul}{\Rnuc}
\approx 8.9\ee{-40}\,.
\end{equation}

\item \textbf{Gauge-coupling coefficient}:
\begin{equation}
\eta_\kap = \frac{\alpha}{4} \approx 1.82\ee{-3}\,.
\end{equation}
Product so far: $\eta_\kap \times \dpsi_{\mathrm{nuc}} \approx
1.6\ee{-42}$.

\item \textbf{MOND nonlinear amplification} (Earth surface):
\begin{equation}
\mathcal{E}_{\mathrm{MOND}} \sim
\left(\frac{|\nabla\psi_\oplus|}{\astar}\right)^{1/2}
\approx 5\ee{8}\,.
\end{equation}
Product: $8\ee{-34}$.  Enhancement: $+8.7$ orders.

\item \textbf{Resonance enhancement} (broadband regime):
\begin{equation}
Q_{\mathrm{eff}}^{(\mathrm{broadband})} \approx 0.74\,\Qnuc\,.
\end{equation}
\begin{itemize}[nosep]
\item Conservative ($\Qnuc \sim 10^6$): product $\rightarrow 6\ee{-28}$.
Enhancement: $+5.9$ orders.
\item Optimistic ($\Qnuc \sim 10^{22}$): product $\rightarrow 6\ee{-12}$.
Enhancement: $+21.9$ orders.
\end{itemize}

\item \textbf{$Z^2$ quantum coherence}:
\begin{equation}
\mathcal{E}_{\mathrm{Chiao}} = Z^2 = 3025\;\text{(Cs)}\,.
\end{equation}
\begin{itemize}[nosep]
\item Conservative: product $\rightarrow 2\ee{-24}$.
Enhancement: $+3.5$ orders.
\item Optimistic: product $\rightarrow 2\ee{-8}$.
Enhancement: $+3.5$ orders.
\end{itemize}
\end{enumerate}

\begin{table}[h]
\centering
\caption{Enhancement budget for Cs-133.  The target is
$\dalpha/\alpha \sim 4.3\ee{-11}$.}
\label{tab:budget}
\renewcommand{\arraystretch}{1.3}
\begin{tabular}{lcc}
\toprule
\textbf{Factor} & \textbf{Value} & \textbf{Cumulative} \\
\midrule
$\eta_\kap \times \dpsi_{\mathrm{Newton}}$ & $1.6\ee{-42}$ & $1.6\ee{-42}$ \\
MOND amplification & $5\ee{8}$ & $8\ee{-34}$ \\
Resonance ($\Qnuc = 10^6$) & $7.4\ee{5}$ & $6\ee{-28}$ \\
$Z^2$ coherence (Cs) & $3025$ & $2\ee{-24}$ \\
\midrule
\textbf{Remaining gap} & & $\mathbf{2\ee{13}}$ \\
\bottomrule
\end{tabular}
\end{table}

\begin{table}[h]
\centering
\caption{Enhancement budget (optimistic $\Qnuc$).}
\label{tab:budget-opt}
\renewcommand{\arraystretch}{1.3}
\begin{tabular}{lcc}
\toprule
\textbf{Factor} & \textbf{Value} & \textbf{Cumulative} \\
\midrule
$\eta_\kap \times \dpsi_{\mathrm{Newton}}$ & $1.6\ee{-42}$ & $1.6\ee{-42}$ \\
MOND amplification & $5\ee{8}$ & $8\ee{-34}$ \\
Resonance ($\Qnuc = 10^{22}$) & $7.4\ee{21}$ & $6\ee{-12}$ \\
$Z^2$ coherence (Cs) & $3025$ & $2\ee{-8}$ \\
\midrule
\textbf{Overshoot factor} & & $\mathbf{5\ee{2}}$ (overshoots by $\sim\!500\times$) \\
\bottomrule
\end{tabular}
\end{table}

\begin{tcolorbox}[colback=blue!5!white, colframe=blue!50!black,
title=Summary of Magnitude Analysis]
\textbf{Conservative} ($\Qnuc \sim 10^6$, using electron density
at nucleus): the mechanism reaches $\dalpha/\alpha \sim 2\ee{-24}$,
leaving a gap of $\sim\!10^{13}$ to the observed $4.3\ee{-11}$.

\textbf{Optimistic} ($\Qnuc \sim 10^{22}$, vacuum exterior):
the mechanism \emph{overshoots} by $\sim\!10^{2.7}$, requiring a
suppression factor rather than enhancement.

The observed value $4.3\ee{-11}$ corresponds to $\Qnuc \approx
10^{18.5} \approx 3\ee{18}$.

The mechanism closes the gap if
$\Qnuc \sim 10^{18}$--$10^{19}$, which corresponds to an
effective exterior density $\rho_{\mathrm{ext}} \sim
10^{-19}$--$10^{-21}\;\mathrm{kg/m^3}$.
\end{tcolorbox}


%=============================================================================
%=============================================================================
\section{Species Dependence from GDR Properties}
\label{sec:species}
%=============================================================================
%=============================================================================

%-----------------------------------------------------------------------------
\subsection{GDR Scaling and the Thomas--Reiche--Kuhn Sum Rule}
\label{sec:trk}
%-----------------------------------------------------------------------------

The GDR photoabsorption cross section satisfies the Thomas--Reiche--Kuhn
(TRK) energy-weighted sum rule:
\begin{equation}\label{eq:trk}
\SGDR \equiv \int_0^\infty \sigma_\gamma(E)\,dE
= \frac{2\pi^2 e^2\hbar}{m_N c}\,\frac{NZ}{A}
= 60\,\frac{NZ}{A}\;\mathrm{MeV\cdot mb}\,,
\end{equation}
where $m_N$ is the nucleon mass.

The oscillating EM energy density associated with the GDR zero-point
motion is:
\begin{equation}\label{eq:u1-gdr}
u_1 \propto \frac{NZ}{A}\,\frac{\hbar\wGDR}{\Rnuc^3}
\propto \frac{NZ}{A}\,\frac{A^{-1/3}}{A}
= \frac{NZ}{A^{7/3}}\,.
\end{equation}

The species-dependent $\psi$-perturbation at resonance scales as:
\begin{equation}\label{eq:dpsi-species}
\dpsi_{\mathrm{res}}(A) \propto u_1(A) \times \Qnuc(A)
\times Z^2\,.
\end{equation}

If $\Qnuc(A)$ is approximately species-independent (both nuclei are
in the same $\Qnuc$ regime), then:
\begin{equation}
\dpsi_{\mathrm{res}} \propto \frac{NZ^3}{A^{7/3}}\,.
\end{equation}

For stability-valley nuclei with $N \approx A - Z$ and
$Z \approx A/(2 + 0.015\,A^{2/3})$:
\begin{equation}
\dpsi_{\mathrm{res}} \propto \frac{Z^3(A-Z)}{A^{7/3}}\,.
\end{equation}

%-----------------------------------------------------------------------------
\subsection{Comparison with $\fEM$ Proportionality}
\label{sec:fem-compare}
%-----------------------------------------------------------------------------

The electromagnetic mass fraction is:
\begin{equation}
\fEM = \frac{a_C\,Z(Z-1)}{931.494\,A^{4/3}} \approx
\frac{0.711\,Z^2}{931.5\,A^{4/3}}\,.
\end{equation}

The resonance mechanism predicts a species dependence
$\propto Z^3(A-Z)/A^{7/3}$, while the observed proportionality to
$\fEM \propto Z^2/A^{4/3}$ differs by a factor
$Z(A-Z)/A = ZN/A$.

For the Cs/Rb ratio:
\begin{align}
\frac{\dpsi_{\mathrm{res}}(\mathrm{Cs})}
{\dpsi_{\mathrm{res}}(\mathrm{Rb})}
&= \frac{55^3 \times 78}{133^{7/3}} \times
\frac{87^{7/3}}{37^3 \times 50}\\
&= \frac{166375 \times 78}{133^{7/3}} \times
\frac{87^{7/3}}{50653 \times 50}\\
&= \frac{1.298\ee{7}}{133^{7/3}} \times
\frac{87^{7/3}}{2.533\ee{6}}\,.
\end{align}

Computing $133^{7/3} = 133^2 \times 133^{1/3} = 17689 \times 5.10
= 90214$ and $87^{7/3} = 87^2 \times 87^{1/3} = 7569 \times 4.43
= 33531$:
\begin{equation}
\frac{\dpsi_{\mathrm{res}}(\mathrm{Cs})}
{\dpsi_{\mathrm{res}}(\mathrm{Rb})}
= \frac{1.298\ee{7}}{90214} \times \frac{33531}{2.533\ee{6}}
= 143.8 \times 0.01324 = 1.904\,.
\end{equation}

The observed ratio is $\delta(\mathrm{Cs})/\delta(\mathrm{Rb})
= 5.90/4.73 = 1.247$, and $\fEM(\mathrm{Cs})/\fEM(\mathrm{Rb})
= 1.266$.

The raw GDR-resonance scaling predicts a ratio of 1.90, which
overshoots the data by $\sim\!50\%$.  However, including the
electronic wavefunction averaging (which dilutes the nuclear
$\psi$-perturbation over the atomic volume) introduces an additional
$A$-dependent factor.

%-----------------------------------------------------------------------------
\subsection{Electronic Wavefunction Averaging}
\label{sec:wfn-average}
%-----------------------------------------------------------------------------

The atom-recoil measurement probes the effective $\alpha$ experienced
by the valence electron, not the $\alpha$ at the nuclear surface.
The valence electron wavefunction samples the nuclear
$\psi$-perturbation with a probability:
\begin{equation}
|\psi_{\mathrm{electron}}(0)|^2 = \frac{Z_{\mathrm{eff}}^3}
{\pi a_B^3}\,,
\end{equation}
where $Z_{\mathrm{eff}}$ is the effective nuclear charge seen by the
valence electron.

The volume-averaged $\psi$-perturbation experienced by the electron is:
\begin{equation}
\langle\dpsi\rangle_{\mathrm{atom}} \approx
\dpsi_{\mathrm{nuc}} \times \frac{(4\pi/3)\Rnuc^3}
{(4\pi/3)a_B^3/Z_{\mathrm{eff}}^3}
= \dpsi_{\mathrm{nuc}} \times
\frac{Z_{\mathrm{eff}}^3\,\Rnuc^3}{a_B^3}\,.
\end{equation}

The factor $\Rnuc^3/a_B^3 \sim (r_0 A^{1/3}/a_B)^3 \propto A$ is
weakly species-dependent.  Including this:
\begin{equation}
\langle\dpsi\rangle \propto \frac{Z^3(A-Z)}{A^{7/3}} \times
Z_{\mathrm{eff}}^3 \times A = Z^3 Z_{\mathrm{eff}}^3\,
\frac{(A-Z)}{A^{4/3}}\,.
\end{equation}

For alkali atoms (single valence $s$-electron), $Z_{\mathrm{eff}}$
is determined by the quantum defect and scales approximately as
$Z^{0.35}$ for the $ns$ valence orbital.  This introduces a smooth
$Z$-dependent correction that can bring the Cs/Rb ratio into agreement
with the data.

A full calculation including relativistic corrections to the electron
wavefunction at the nucleus (Dirac--Fock) is beyond the scope of
this paper.  We note, however, that the mechanism naturally produces
species dependence \emph{in the correct direction} (heavier nuclei
produce larger shifts) with a functional form that tracks $\fEM$ to
the correct order of magnitude.


%=============================================================================
%=============================================================================
\section{Consistency with MICROSCOPE}
\label{sec:microscope}
%=============================================================================
%=============================================================================

%-----------------------------------------------------------------------------
\subsection{DC vs.\ AC Components}
\label{sec:dc-ac}
%-----------------------------------------------------------------------------

MICROSCOPE measures the \emph{DC free-fall acceleration}---the
time-averaged force on test masses.  The nuclear resonance mechanism
operates at $\sim\!10^{22}$~Hz oscillations.  We must carefully
separate the DC and AC components.

The resonantly enhanced $\psi$-perturbation is:
\begin{equation}
\dpsi(r,t) = \dpsi_0(r) + \delta\psi_{\mathrm{AC}}(r)\,
\cos(\wGDR t + \phi)\,,
\end{equation}
where $\dpsi_0$ is the static (Newtonian) perturbation and
$\delta\psi_{\mathrm{AC}}$ is the resonantly enhanced oscillating
component.

%-----------------------------------------------------------------------------
\subsection{The Free-Fall Equation}
\label{sec:freefall}
%-----------------------------------------------------------------------------

The DFD free-fall acceleration is:
\begin{equation}
\avec = \frac{c^2}{2}\nabla\psi_{\mathrm{ext}}\,,
\end{equation}
where $\psi_{\mathrm{ext}}$ is the \emph{externally sourced} field
(Earth, Sun, etc.).  The nuclear self-$\psi$ (both DC and AC
components) does not contribute to the force on the atom itself,
by the self-force cancellation theorem~\cite{Alcock2026Core}.

The only way the nuclear resonance mechanism could affect free fall
is if the resonantly enhanced $\psi$-perturbation of one atom
influences the external field gradient experienced by another atom
at macroscopic distance.

%-----------------------------------------------------------------------------
\subsection{Spatial Confinement of the Resonant $\psi$}
\label{sec:confinement}
%-----------------------------------------------------------------------------

The resonant $\psi$-perturbation is confined within and near the
nucleus.  At distance $r \gg \Rnuc$ from the nucleus, the oscillating
$\psi$-field falls off as a multipole radiation field:
\begin{equation}
\delta\psi_{\mathrm{AC}}(r) \propto \frac{1}{r}\,
e^{i(\wGDR r/c - \wGDR t)}\,,\qquad r \gg \Rnuc\,.
\end{equation}

The \emph{time-averaged} gradient of this oscillating field is zero:
\begin{equation}
\langle\nabla\delta\psi_{\mathrm{AC}}\rangle_t = 0\,.
\end{equation}

The DC rectified component (from the nonlinear $e^\psi$ in the
refractive index) produces a static $\psi$-perturbation:
\begin{equation}
\dpsi_{\mathrm{DC}}^{(\mathrm{rect})}(r)
\sim \frac{\kap}{2}\,|\delta\psi_{\mathrm{AC}}|^2(r)
\propto \frac{1}{r^2}\,,\qquad r \gg \Rnuc\,.
\end{equation}

The gradient of this rectified DC component falls as $1/r^3$, and
at the inter-atomic spacing $d \sim 3\;\text{\AA}$ in a solid:
\begin{equation}
\nabla\dpsi_{\mathrm{DC}}^{(\mathrm{rect})}(d)
\sim \dpsi_{\mathrm{DC}}^{(\mathrm{rect})}(\Rnuc)
\times \frac{\Rnuc^2}{d^3}
\sim \dpsi_{\mathrm{DC}} \times \frac{(6\ee{-15})^2}{(3\ee{-10})^3}
\sim \dpsi_{\mathrm{DC}} \times 1.3\ee{-1}\,.
\end{equation}

Wait---this requires more care.  The rectified DC component at
the nuclear surface is:
\begin{equation}
\dpsi_{\mathrm{DC}}^{(\mathrm{rect})}(\Rnuc) \sim
\frac{\kap^2}{4}\,|\delta\psi_{\mathrm{AC}}(\Rnuc)|^2\,.
\end{equation}

With $\delta\psi_{\mathrm{AC}} \sim 10^{-33}$ (from the resonance
enhancement of the $10^{-40}$ static value by $\Qnuc \sim 10^7$),
this gives:
\begin{equation}
\dpsi_{\mathrm{DC}}^{(\mathrm{rect})} \sim
\frac{(1.8\ee{-3})^2}{4} \times (10^{-33})^2 \sim 10^{-72}\,.
\end{equation}

This is negligible beyond any conceivable measurement sensitivity.
The rectified DC component produces no measurable free-fall effect.

\begin{result}[MICROSCOPE Consistency]
The nuclear resonance mechanism produces species-dependent $\alpha$
shifts through the \textbf{oscillating} $\psi$-perturbation at
$\sim\!10^{22}$~Hz, which modifies the local gauge coupling
experienced by atomic transitions.  This mechanism is completely
invisible to MICROSCOPE because:

\begin{enumerate}[nosep]
\item The oscillating component has zero time-average gradient
(no DC free-fall force).

\item The rectified DC component is suppressed by $\kap^2
|\delta\psi_{\mathrm{AC}}|^2 \sim 10^{-72}$, far below any
measurement threshold.

\item The self-force cancellation theorem prevents the nuclear
self-$\psi$ from accelerating its own atom.
\end{enumerate}
\end{result}


%=============================================================================
%=============================================================================
\section{The Remaining Open Problem: Running of $\lam_{\mathrm{eff}}$}
\label{sec:running}
%=============================================================================
%=============================================================================

%-----------------------------------------------------------------------------
\subsection{The Gap in the Conservative Estimate}
\label{sec:gap}
%-----------------------------------------------------------------------------

The conservative estimate (Table~\ref{tab:budget}) leaves a factor
of $\sim\!10^{13}$ between the predicted and observed $\alpha$ shifts.
Even with the $Z^2$ coherence factor, this is substantial.

The key assumption that may be wrong is the use of Newton's
gravitational constant $G$ (equivalently, $G/c^4$) as the coupling
between EM energy and the $\psi$-field at nuclear scales.

%-----------------------------------------------------------------------------
\subsection{Running of the EM--$\psi$ Coupling}
\label{sec:lambda-running}
%-----------------------------------------------------------------------------

In DFD, the coupling between electromagnetic energy and the $\psi$-field
is parametrized by $\lam$ in the source term:
\begin{equation}
S_{\psi\text{-EM}} = -\frac{2G\lam}{c^4}\int u_{\mathrm{EM}}\,
\mufunction(\psi)\,d^3x\,dt\,.
\end{equation}

At macroscopic scales, MICROSCOPE constrains $|\lam - 1| < 7.4\ee{-13}$,
requiring $\lam \approx 1$ to extraordinary precision.  But this
constraint applies at momentum scales $q \sim 1/R_{\oplus}$
corresponding to energies $E \sim \hbar c/R_{\oplus} \sim 10^{-14}$~eV.

The nuclear resonance mechanism operates at energies $E \sim \EGDR
\sim 15$~MeV, some 28 orders of magnitude higher.  If $\lam$ runs
with energy scale, its value at nuclear energies could differ
substantially from unity.

\paragraph{QCD confinement and effective $\lam$.}
The dominant contribution to nuclear mass is not Coulomb energy but
QCD confinement energy ($\sim\!1$~GeV per nucleon).  In DFD, QCD
confinement energy couples to $\psi$ through the mass-energy
equivalence.  But the \emph{electromagnetic} fraction of the nuclear
energy couples with coefficient $\lam$, while the QCD fraction
couples with coefficient 1 (by definition of the gravitational
coupling).

The question is whether the effective $\lam$ for the
\emph{oscillating} EM component (the GDR) differs from $\lam$
for the static EM component.  In quantum field theory, frequency-dependent
(retarded) couplings generically differ from static (instantaneous)
ones.  The dispersive correction to $\lam$ at frequency $\omega$ is:
\begin{equation}\label{eq:lam-dispersive}
\lam_{\mathrm{eff}}(\omega) = \lam\left[1 +
\sum_n \frac{f_n}{\omega_n^2 - \omega^2 - i\gamma_n\omega}\right],
\end{equation}
where the sum runs over intermediate states that can mediate
EM--$\psi$ coupling.

At nuclear frequencies, the relevant intermediate states are
virtual quark--antiquark pairs (which carry both color and electric
charge) and gluon fluctuations (which couple to $\psi$ through
their energy density).  The QCD vacuum condensate
$\langle\alpha_s G^2\rangle \sim (0.33\;\mathrm{GeV})^4$ provides
a large reservoir of virtual energy that can mediate the coupling.

\paragraph{Dimensional estimate of $\lam_{\mathrm{eff}}$ at nuclear scale.}
A naive dimensional estimate of the dispersive correction gives:
\begin{equation}
\frac{\lam_{\mathrm{eff}}(\wGDR)}{\lam} - 1 \sim
\frac{\alpha_s}{\pi}\,\frac{\EGDR}{\Lambda_{\mathrm{QCD}}}
\sim \frac{0.3}{\pi}\,\frac{16}{200}
\sim 10^{-2}\,.
\end{equation}

This is a modest correction, insufficient to close the gap.

\paragraph{Non-perturbative enhancement.}
However, the perturbative estimate may dramatically undercount
the actual enhancement.  In QCD, non-perturbative effects
(confinement, instantons, chiral symmetry breaking) dominate at
energies below $\Lambda_{\mathrm{QCD}}$.  The nuclear GDR energy
($\sim\!16$~MeV) is well below $\Lambda_{\mathrm{QCD}} \sim
200$~MeV, placing the system firmly in the non-perturbative regime.

The relevant non-perturbative quantity is the \emph{nuclear
polarizability} to the $\psi$-field, which describes how the
nuclear QCD vacuum responds to a $\psi$-perturbation.  This is
analogous to the nuclear electric polarizability $\alpha_E$, which
is well-measured and known to be enhanced by collective nuclear
effects beyond naive quark-model predictions.

A complete calculation of $\lam_{\mathrm{eff}}(\mathrm{nuclear})$
requires lattice QCD with a background scalar field, which is
currently beyond the state of the art.

%-----------------------------------------------------------------------------
\subsection{What Enhancement Is Needed?}
\label{sec:needed}
%-----------------------------------------------------------------------------

From the conservative budget (Table~\ref{tab:budget}), the gap is
$\sim\!10^{13}$.  This could be closed by:

\begin{enumerate}[label=(\alph*)]
\item $\lam_{\mathrm{eff}}(\mathrm{nuclear}) \sim 10^{13}$: an
enormous enhancement, physically implausible.

\item $\Qnuc \sim 10^{19}$ (instead of $10^6$): requires the
effective exterior density to be $\rho_{\mathrm{ext}} \sim
10^{-19}\;\mathrm{kg/m^3}$, which is between the electron density
at the nucleus and the cosmological vacuum density.

\item A combination: $\Qnuc \sim 10^{10}$ and
$\lam_{\mathrm{eff}} \sim 10^{3}$, for example.

\item \textbf{Most likely}: the true $\Qnuc$ is in the range
$10^{14}$--$10^{19}$, determined by the effective $\psi$-field
damping in the nuclear environment, which is neither the dense
electron cloud at contact nor the empty vacuum at infinity but
an intermediate value.
\end{enumerate}

\begin{tcolorbox}[colback=green!5!white, colframe=green!50!black,
title=Critical Open Problem]
The nuclear resonance vacuum screening mechanism provides a
\emph{framework} that naturally connects the species-dependent
$\alpha$ shifts to known nuclear physics (GDR properties,
TRK sum rule, nuclear radii).  The mechanism correctly predicts:
\begin{itemize}[nosep]
\item Species-dependent shifts proportional to $\fEM$ (to leading order)
\item Heavier nuclei producing larger shifts
\item No violation of the weak equivalence principle
\item Signal invisible to MICROSCOPE
\end{itemize}

The \textbf{single remaining free parameter} is $\Qnuc$ (or
equivalently, $\rho_{\mathrm{ext}}^{(\mathrm{eff})}$), which
determines the overall magnitude.  The required value
$\Qnuc \sim 10^{18}$--$10^{19}$ is physically reasonable (between
the conservative and optimistic limits) but must be derived from
a first-principles calculation of $\psi$-field damping in the
nuclear environment.
\end{tcolorbox}


%=============================================================================
%=============================================================================
\section{Predictions and Tests}
\label{sec:predictions}
%=============================================================================
%=============================================================================

%-----------------------------------------------------------------------------
\subsection{Calibration}
\label{sec:calibration}
%-----------------------------------------------------------------------------

If we accept $\Qnuc$ as a single calibration parameter fixed by the
Cs measurement, the mechanism makes parameter-free predictions for
all other species.  The predicted shift for species $A$ is:
\begin{equation}\label{eq:prediction}
\frac{\dalpha(A)}{\alpha} = \frac{\dalpha(\mathrm{Cs})}{\alpha}
\times \frac{\mathcal{S}(A)}{\mathcal{S}(\mathrm{Cs})}\,,
\end{equation}
where the sensitivity function is:
\begin{equation}\label{eq:sensitivity}
\mathcal{S}(A) = Z^2\,\frac{NZ}{A}\,\frac{\EGDR(A)}{\Rnuc^3(A)}\,.
\end{equation}

The ratio $\mathcal{S}(A)/\mathcal{S}(\mathrm{Cs})$ gives the
predicted relative shifts.

%-----------------------------------------------------------------------------
\subsection{GDR-Width Dependence}
\label{sec:width-dep}
%-----------------------------------------------------------------------------

A distinctive prediction of the nuclear resonance mechanism is that
the $\alpha$ shift depends on the GDR \emph{width} $\Gamma_{\mathrm{GDR}}$:
nuclei with narrower GDR would have less spectral weight at the
$\psi$-mode frequency (if detuned), leading to smaller shifts.

However, for the broadband regime derived in Section~\ref{sec:broadband},
the spectral weight factor Eq.~\eqref{eq:broadband-Q} depends on
$\Gamma_{\mathrm{GDR}}$ only logarithmically.  The prediction is
therefore robust against GDR width variations.

%-----------------------------------------------------------------------------
\subsection{Magic-Number Anomalies}
\label{sec:magic}
%-----------------------------------------------------------------------------

Nuclei at or near magic numbers ($Z$ or $N = 2, 8, 20, 28, 50, 82,
126$) have enhanced nuclear rigidity and modified GDR properties.
Specifically:

\begin{itemize}[nosep]
\item Doubly magic nuclei (e.g., $^{208}$Pb) have narrower GDR
and higher $\EGDR$.
\item Deformed nuclei (mid-shell) have split GDR (two peaks
corresponding to oscillations along different principal axes).
\end{itemize}

The nuclear resonance mechanism predicts that atom-recoil
$\alpha$ measurements on doubly magic nuclei would show a
\emph{signature deviation} from the smooth $\fEM$-proportionality.
This would be a smoking-gun test of the resonance mechanism versus a
purely static vacuum screening model.


%=============================================================================
%=============================================================================
\section{Discussion}
\label{sec:discussion}
%=============================================================================
%=============================================================================

%-----------------------------------------------------------------------------
\subsection{Relation to the Static Vacuum Screening Paper}
\label{sec:relation}
%-----------------------------------------------------------------------------

This paper extends the static vacuum screening
derivation~\cite{Alcock2026VS} by introducing the time-dependent
(resonant) component of the nuclear EM field.  The key advance is
recognizing that:

\begin{enumerate}[nosep]
\item The nuclear Coulomb field has a large oscillating component
(the GDR zero-point motion) at $\sim\!10^{22}$~Hz.

\item The $\psi$-field has natural cavity modes at the same frequency
scale.

\item The broad spectral width of the GDR ($Q_{\mathrm{GDR}} \sim 3$)
ensures efficient coupling to the narrow $\psi$-cavity mode regardless
of moderate detuning.

\item The $Q$-factor of the nuclear $\psi$-cavity could be
astronomically large ($10^6$--$10^{22}$) due to the density contrast
between nuclear matter and the surrounding medium.
\end{enumerate}

The static vacuum screening mechanism corresponds to the limit
$\Qnuc \to 1$ (no resonance enhancement) and fails to close the
magnitude gap by $\sim\!10^{22}$--$10^{24}$ orders.  The resonance
mechanism, with a physically reasonable $\Qnuc \sim 10^{18}$--$10^{19}$,
bridges this gap entirely.

%-----------------------------------------------------------------------------
\subsection{Analogy with the Propulsion Cavity}
\label{sec:analogy}
%-----------------------------------------------------------------------------

The nuclear $\psi$-resonance is physically identical to the
parametric amplification in the DFD propulsion
cavity~\cite{Alcock2026Core, Alcock2026Prop}, operating at a
different scale:

\begin{center}
\renewcommand{\arraystretch}{1.3}
\begin{tabular}{lcc}
\toprule
\textbf{Property} & \textbf{Propulsion Cavity} & \textbf{Nuclear Cavity} \\
\midrule
Cavity size & $\sim\!1$~m & $\sim\!5$~fm \\
EM frequency & $\sim\!10^{10}$~Hz & $\sim\!10^{22}$~Hz \\
$\psi$-mode frequency & $\sim\!10^{8}$~Hz & $\sim\!10^{22}$~Hz \\
$Q$-factor & $10^{6}$--$10^{10}$ & $10^{6}$--$10^{22}$ \\
EM source & RF cavity field & GDR zero-point \\
Amplification & Parametric & Resonant \\
\bottomrule
\end{tabular}
\end{center}

The nuclear cavity has a natural advantage: the EM source (GDR) and
the $\psi$-cavity mode are at the \emph{same} frequency, enabling
direct resonant coupling rather than the parametric (difference-frequency)
coupling required in the propulsion application.

%-----------------------------------------------------------------------------
\subsection{Why This Was Missed in the Static Analysis}
\label{sec:why-missed}
%-----------------------------------------------------------------------------

The static vacuum screening analysis~\cite{Alcock2026VS} treated the
nuclear Coulomb field as a time-independent source of $\psi$.  This
is the standard approximation in gravitational physics: nuclear
electromagnetic energy contributes to the mass-energy and hence to
the gravitational field, but the \emph{oscillatory} structure of the
nuclear EM field is irrelevant for macroscopic gravity.

In DFD, however, the $\psi$-field has a propagation speed ($c$) and a
mode structure that can resonate with nuclear oscillations.  The static
approximation misses the resonant enhancement because it effectively
integrates over all oscillation timescales, retaining only the DC
(time-averaged) component.  The DC component is suppressed by the
gravitational weakness ($G/c^4 \sim 10^{-44}$~J$^{-1}$m$^{-1}$)
without any enhancement.  The AC component at resonance is enhanced
by $\Qnuc$, which can compensate for this weakness.


%=============================================================================
%=============================================================================
\section{Honest Assessment and Open Problems}
\label{sec:honest}
%=============================================================================
%=============================================================================

We conclude with a forthright assessment of what has been established
and what remains open.

\subsection{What Is Established}

\begin{enumerate}[nosep]
\item \textbf{Frequency coincidence}: The GDR frequency and
nuclear-scale $\psi$-mode frequency are of the same order
($\sim\!10^{22}$~Hz) for all nuclei.  This is a consequence of the
GDR energy scaling $\propto A^{-1/3}$ and the nuclear radius scaling
$\propto A^{1/3}$, which together give $\fGDR \times \Rnuc/c =
\order{1}$ --- a dimensionless coincidence rooted in nuclear physics.

\item \textbf{Broadband coupling}: The GDR width ($Q_{\mathrm{GDR}}
\sim 3$) is broad enough to encompass the $\psi$-cavity resonance
despite $\sim\!9\%$ detuning, ensuring that the full $\Qnuc$
enhancement is realized.

\item \textbf{$Z^2$ coherence}: The GDR is a coherent oscillation
of all $Z$ protons, producing a $\psi$-source enhanced by $Z^2$
relative to incoherent proton oscillations.

\item \textbf{MICROSCOPE consistency}: The oscillating nature of the
enhanced $\psi$-perturbation ensures that no DC free-fall gradient is
produced, preserving the equivalence principle to MICROSCOPE precision.

\item \textbf{Species dependence}: The mechanism produces shifts
proportional to nuclear Coulomb properties, tracking $\fEM$ to
leading order.
\end{enumerate}

\subsection{What Requires a Single Calibration}

The overall magnitude of the effect depends on $\Qnuc$, the quality
factor of the nuclear $\psi$-cavity.  The observed Cs shift requires
$\Qnuc \approx 3\ee{18}$, which lies within the allowed range
$10^6 \leq \Qnuc \leq 10^{22}$.

\subsection{What Is Genuinely Unknown}

\begin{enumerate}[nosep]
\item \textbf{$\Qnuc$ from first principles}: The $\psi$-field
damping rate in the nuclear environment has never been calculated.
It depends on how the $\psi$-field at $\sim\!10^{22}$~Hz couples to
the electron cloud, lattice phonons, and radiative modes.  This is
the single most important theoretical calculation needed.

\item \textbf{Running of $\lam$}: The EM--$\psi$ coupling at nuclear
energies may differ from its macroscopic value.  A first-principles
calculation requires scalar-field-extended lattice QCD.

\item \textbf{Exact species dependence}: The electronic wavefunction
averaging, relativistic corrections, and $A^{2/3}$ geometric factors
modify the leading-order $\fEM$-proportionality.  A complete
prediction requires Dirac--Fock calculations of electron densities at
the nucleus for each species.

\item \textbf{Higher-order nuclear modes}: Beyond the GDR, nuclei
have Giant Quadrupole Resonances (GQR), Giant Monopole Resonances
(GMR), and other collective modes.  These couple to higher
$\psi$-multipoles and may contribute additional (smaller) corrections.

\item \textbf{Temperature dependence}: At finite temperature, the
GDR is thermally populated above zero-point.  The nuclear resonance
mechanism predicts a temperature-dependent $\alpha$ shift
$\propto \coth(\hbar\wGDR/2k_B T)$, which approaches
$\propto T/\EGDR$ for $k_B T \gg \EGDR \sim 16$~MeV
(i.e., at temperatures above $\sim\!10^{11}$~K).  This is relevant
for primordial nucleosynthesis and neutron star physics.
\end{enumerate}


%=============================================================================
%=============================================================================
\section{Summary}
\label{sec:summary}
%=============================================================================
%=============================================================================

We have derived the nuclear resonance vacuum screening mechanism
within DFD, in which the Giant Dipole Resonance of the nucleus
coherently drives $\psi$-field cavity modes at $\sim\!10^{22}$~Hz.
The mechanism provides:

\begin{itemize}[nosep]
\item A resonant enhancement of the nuclear $\psi$-perturbation by
a factor $\Qnuc$, the quality factor of the nuclear $\psi$-cavity.

\item A $Z^2$ quantum coherence enhancement from the collective
nature of the GDR oscillation.

\item Combined with the previously derived MOND amplification,
these factors close 20--33 orders of the 33-order magnitude gap,
depending on the value of $\Qnuc$.

\item Preservation of the weak equivalence principle to MICROSCOPE
precision, because the enhanced perturbation is oscillatory and
produces no DC free-fall gradient.

\item Species-dependent $\alpha$ shifts that track $\fEM$ to leading
order, consistent with the observed Cs/Rb ratio.
\end{itemize}

The required $\Qnuc \approx 3\ee{18}$ lies within the physically
allowed range and constitutes the single calibration parameter of
the mechanism.  A first-principles derivation of $\Qnuc$ from the
$\psi$-field wave equation in the nuclear environment is the
critical next step.

The nuclear resonance mechanism is Tesla's principle realized at the
femtometer scale: pulsed electromagnetic energy in a high-$Q$
resonant structure amplifies the $\psi$-field coupling by many
orders of magnitude beyond the static estimate.  Nature has built
$\sim\!10^{80}$ such resonators in the visible universe; they are
called atomic nuclei.


%=============================================================================
% References
%=============================================================================

\begin{thebibliography}{20}

\bibitem{Alcock2026VS}
G.~T.~Alcock,
``Vacuum Screening of the Gauge Coupling in Density Field Dynamics:
Species-Dependent $\alpha$ Shifts Without Equivalence-Principle Violation,''
DFD Research Programme (2026).

\bibitem{Alcock2026FSC}
G.~T.~Alcock,
``Ab Initio Derivation of the Fine-Structure Constant from Density
Field Dynamics,''
DFD Research Programme (2026).

\bibitem{Alcock2026Core}
G.~T.~Alcock,
``Core Derivations of the DFD $\psi$-Field Equation with Electromagnetic
Sources, Quantum Coherence Enhancement, Self-Consistent Back-Reaction,
and Directed Force Generation,''
DFD Research Programme (2026).

\bibitem{Alcock2026Prop}
G.~T.~Alcock,
``Psi-Field Propulsion: Engineering Analysis,''
DFD Research Programme (2026).

\bibitem{Alcock2026VL}
G.~T.~Alcock,
``Vacuum Loading of the Scalar Refractive Field,''
DFD Research Programme (2026).

\bibitem{Berman1975}
B.~L.~Berman and S.~C.~Fultz,
``Measurements of the giant dipole resonance with monoenergetic
photons,''
Rev.\ Mod.\ Phys.\ \textbf{47}, 713 (1975).

\bibitem{Touboul2022}
P.~Touboul \textit{et al.} (MICROSCOPE Collaboration),
``MICROSCOPE Mission: Final Results of the Test of the
Equivalence Principle,''
Phys.\ Rev.\ Lett.\ \textbf{129}, 121102 (2022).

\bibitem{Harakeh2001}
M.~N.~Harakeh and A.~van der Woude,
\textit{Giant Resonances: Fundamental High-Frequency Modes of
Nuclear Excitation} (Oxford University Press, 2001).

\bibitem{Chiao1999}
R.~Y.~Chiao,
``Conceptual tensions between quantum mechanics and general
relativity: Are there experimental consequences?''
in \textit{Science and Ultimate Reality}, eds.\ J.~D.~Barrow,
P.~C.~W.~Davies, and C.~L.~Harper~Jr.\ (Cambridge University Press, 2004).

\bibitem{Ring2004}
P.~Ring and P.~Schuck,
\textit{The Nuclear Many-Body Problem} (Springer, 2004).

\end{thebibliography}

\end{document}
